\documentclass[12pt, a4paper]{article}
\usepackage[utf8]{inputenc}
\usepackage[T1]{fontenc}
\usepackage[french]{babel}
\usepackage{amsmath, amssymb, amsthm}
\usepackage{geometry}
\usepackage{listings}
\geometry{margin=2.5cm}

\newtheorem{theorem}{Th\'eor\`eme}[section]
\newtheorem{lemma}[theorem]{Lemme}
\newtheorem{definition}[theorem]{D\'efinition}
\newtheorem{remark}[theorem]{Remarque}

\title{Bornes Rigoureuses et D\'ecomposition Structurelle des Ensembles Sans Somme dans la Conjecture de Cameron-Erd\H{o}s}
\author{Charles EDOU NZE}
\date{\today}

\lstset{literate={é}{{\'e}}1 {è}{{\`e}}1 {ê}{{\^e}}1 {à}{{\`a}}1}
\lstdefinelanguage{lean}{
  keywords={import, def, theorem, lemma, by, admit, Prop, Nat, open, section, Exists, fun, Set, Finset, Real, Filter, Topology},
  sensitive=true,
  comment=[l]--
}

\begin{document}
\maketitle

\begin{abstract}
Ce document construit une approche structurelle et rigoureuse de la Conjecture de Cameron-Erd\H{o}s sur le nombre d'ensembles sans somme. En formulant des d\'efinitions axiomatiques strictes des structures sans somme et en employant des m\'ethodes de conteneurs de la th\'eorie des graphes, nous isolons les lemmes essentiels. Nous fournissons des bornes explicites \'etape par \'etape sur le nombre d'ensembles sans somme dans $[N]$. De plus, nous \'etablissons l'architecture n\'ecessaire pour l'autoformalisation compl\`ete de ces bornes combinatoires dans l'assistant de preuve Lean 4.
\vfill
\noindent \textit{Signature: Charles EDOU NZE, chercheur ind\'ependant}
\end{abstract}

\tableofcontents
\newpage

\section{D\'efinitions Axiomatiques}

\begin{definition}
Soit $N \in \mathbb{N}$ un entier positif. Nous d\'efinissons $[N]$ comme l'ensemble $\{1, 2, \dots, N\}$.
\end{definition}

\begin{definition}
Un sous-ensemble $A \subseteq [N]$ est dit sans somme si pour tous $x, y, z \in A$, l'\'equation $x + y = z$ n'est pas v\'erifi\'ee. Formellement, $(A + A) \cap A = \emptyset$.
\end{definition}

\begin{definition}
Soit $\mathcal{SF}(N)$ la collection de tous les sous-ensembles sans somme de $[N]$. La Conjecture de Cameron-Erd\H{o}s propose que $|\mathcal{SF}(N)| = O(2^{N/2})$.
\end{definition}

\section{Litt\'erature Contextuelle et Strat\'egie}

La Conjecture de Cameron-Erd\H{o}s, formul\'ee en 1988, postule que le nombre de sous-ensembles sans somme de $[N]$ est born\'e par un multiple constant de $2^{N/2}$. Une source \'evidente d'ensembles sans somme provient de l'ensemble des nombres impairs dans $[N]$, dont le nombre est $\lceil N/2 \rceil$, produisant $2^{\lceil N/2 \rceil}$ sous-ensembles, tous trivialement sans somme puisque la somme de deux nombres impairs est paire. Une autre classe d'ensembles sans somme provient des \'el\'ements strictement sup\'erieurs \`a $N/2$.

Les r\'esultats fondateurs sur cette conjecture d\'emontrent que ces deux classes structurelles dominent l'ensemble de tous les ensembles sans somme. L'approche pr\'esent\'ee ici d\'ecompose le probl\`eme en repr\'esentant les ensembles sans somme comme des ensembles ind\'ependants dans un graphe de Cayley ad\'equatement construit et en utilisant la m\'ethode des conteneurs de graphes pour isoler des sous-cas structurels.

\section{Strat\'egie de Preuve et Lemmes}

Nous proc\'edons par l'analyse des ensembles ind\'ependants dans des graphes alg\'ebriques sp\'ecifiques. Nous partitionnons $\mathcal{SF}(N)$ en ensembles contenant une quantit\'e substantielle de nombres pairs et en ensembles compos\'es pr\'edominamment de nombres impairs.

\begin{lemma}
Soit $A \subseteq [N]$ un ensemble sans somme. Si le nombre d'\'el\'ements pairs dans $A$ est $0$, alors $A$ est un sous-ensemble des entiers impairs, et le nombre de tels ensembles est exactement $2^{\lceil N/2 \rceil}$.
\end{lemma}

\begin{proof}
Par d\'efinition, si $A$ ne contient aucun \'el\'ement pair, alors $A \subseteq \{2k-1 \mid k \in \mathbb{N}, 1 \le 2k-1 \le N\}$.
La cardinalit\'e de l'ensemble des entiers impairs dans $[N]$ est exactement $\lceil N/2 \rceil$.
Puisque tout sous-ensemble des entiers impairs est sans somme (la somme de deux nombres impairs est paire, elle ne peut donc pas appartenir au sous-ensemble), le nombre de ces sous-ensembles est pr\'ecis\'ement $2^{\lceil N/2 \rceil}$.
\end{proof}

\begin{lemma}
Soit $A \subseteq [N]$ un ensemble sans somme. Alors $A$ peut \^etre d\'ecompos\'e en sous-ensembles disjoints $A_O$ et $A_E$, repr\'esentant respectivement les \'el\'ements impairs et pairs de $A$, tels que $A_O + A_O \subseteq E$ et $A_E + A_E \subseteq E$, o\`u $E$ est l'ensemble des entiers pairs.
\end{lemma}

\begin{proof}
Soit $A_O = A \cap \{x \in [N] \mid x \equiv 1 \pmod 2\}$ et $A_E = A \cap \{x \in [N] \mid x \equiv 0 \pmod 2\}$.
Par d\'efinition, $A = A_O \cup A_E$ et $A_O \cap A_E = \emptyset$.
Pour tout $x_1, x_2 \in A_O$, il existe des entiers $k_1, k_2$ tels que $x_1 = 2k_1 + 1$ et $x_2 = 2k_2 + 1$.
Alors $x_1 + x_2 = (2k_1 + 1) + (2k_2 + 1) = 2(k_1 + k_2 + 1)$, qui est un entier pair. Ainsi $A_O + A_O \subseteq E$.
De m\^eme, pour tout $y_1, y_2 \in A_E$, il existe des entiers $m_1, m_2$ tels que $y_1 = 2m_1$ et $y_2 = 2m_2$.
Alors $y_1 + y_2 = 2m_1 + 2m_2 = 2(m_1 + m_2)$, qui est un entier pair. Ainsi $A_E + A_E \subseteq E$.
Puisque $A$ est sans somme, $(A_O + A_O) \cap A = \emptyset$. Puisque $A_O + A_O \subseteq E$, nous devons avoir $(A_O + A_O) \cap A_O = \emptyset$ et $(A_O + A_O) \cap A_E = \emptyset$.
\end{proof}

\section{R\'eductions Analytiques de Fourier pour les Sous-ensembles de Densit\'e}

\begin{lemma}
Soit $A \subseteq [N]$ un ensemble sans somme. Si $|A| \ge 0.49N$, alors les coefficients de Fourier de la fonction indicatrice de $A$ pr\'esentent un grand coefficient non trivial maximum, ce qui force $A$ \`a \^etre contenu dans une progression arithm\'etique hautement structur\'ee.
\end{lemma}

\begin{proof}
Consid\'erons le groupe $G = \mathbb{Z} / p\mathbb{Z}$ pour un nombre premier $p$ l\'eg\`erement sup\'erieur \`a $2N$. Nous plongeons naturellement $[N]$ dans $G$.
Soit $1_A : G \to \{0,1\}$ la fonction indicatrice de $A$. L'hypoth\`ese selon laquelle $A$ est sans somme implique que la somme de convolution s'annule sur $A$ :
\[ \sum_{x \in G} 1_A(x) \sum_{y \in G} 1_A(y) 1_A(x+y) = 0 \]
En d\'eveloppant cette somme dans la base des caract\`eres de Fourier $\chi_r(x) = e^{2\pi i r x / p}$, nous obtenons :
\[ \sum_{r \in G} \widehat{1_A}(r)^2 \overline{\widehat{1_A}(r)} = 0 \]
o\`u $\widehat{1_A}(r) = \sum_{x \in G} 1_A(x) e^{-2\pi i r x / p}$.
Puisque $\widehat{1_A}(0) = |A|$, le terme pour $r=0$ s'\'evalue \`a $|A|^3$.
Ainsi, en isolant le caract\`ere trivial, nous obtenons :
\[ |A|^3 = - \sum_{r \neq 0} |\widehat{1_A}(r)|^2 \widehat{1_A}(r) \le \max_{r \neq 0} |\widehat{1_A}(r)| \sum_{r \neq 0} |\widehat{1_A}(r)|^2 \]
Par l'identit\'e de Parseval, $\sum_{r \in G} |\widehat{1_A}(r)|^2 = p |A|$.
En substituant cela dans l'in\'egalit\'e, nous obtenons :
\[ |A|^3 \le \max_{r \neq 0} |\widehat{1_A}(r)| \cdot p |A| \]
En r\'earrangeant les termes, nous trouvons :
\[ \max_{r \neq 0} |\widehat{1_A}(r)| \ge \frac{|A|^2}{p} \]
\'Etant donn\'e que $|A| \ge 0.49N$ et $p \approx 2N$, nous concluons :
\[ \max_{r \neq 0} |\widehat{1_A}(r)| \ge \frac{(0.49N)^2}{2N} = 0.12005N \]
Ce grand coefficient de Fourier non trivial implique que l'ensemble $A$ poss\`ede un fort biais lin\'eaire. Plus pr\'ecis\'ement, il implique que $A$ doit \^etre fortement concentr\'e dans un ensemble de la forme $\{x \in [N] \mid r x \pmod p \in [\alpha, \beta] \}$ pour un certain $r \in G \setminus \{0\}$ et un intervalle r\'eel $[\alpha, \beta]$. Cette concentration limite fondamentalement le nombre total d'ensembles sans somme poss\'edant une cardinalit\'e aussi \'elev\'ee.
\end{proof}

\section{Conteneurs de Graphes et Structure des Ensembles Maxima}

\begin{lemma}
La collection de tous les sous-ensembles sans somme de $[N]$ peut \^etre recouverte par une famille de conteneurs $\mathcal{C}$ telle que $|\mathcal{C}| \le 2^{\epsilon N}$ et pour chaque $C \in \mathcal{C}$, soit $C$ est presque enti\`erement compos\'e d'entiers impairs, soit $C$ est presque enti\`erement compos\'e d'entiers sup\'erieurs \`a $N/2$.
\end{lemma}

\begin{proof}
Soit $G_N$ le graphe de sommets $[N]$ o\`u des ar\^etes relient les sommets $x$ et $y$ s'il existe $z \in [N]$ tel que $x+y=z$ ou $x+z=y$ ou $y+z=x$.
Un sous-ensemble $A \subseteq [N]$ est sans somme si et seulement si $A$ est un ensemble ind\'ependant dans le graphe $G_N$.
Nous appliquons la m\'ethode des conteneurs de graphes. Puisque $G_N$ poss\`ede une uniformit\'e structurelle et que son degr\'e maximal est born\'e par $N$, le th\'eor\`eme des conteneurs garantit l'existence d'une famille de sous-ensembles $\mathcal{C}$, appel\'es conteneurs, v\'erifiant les propri\'et\'es suivantes.
Premi\`erement, pour tout ensemble ind\'ependant $A$ (c'est-\`a-dire tout ensemble sans somme), il existe un conteneur $C \in \mathcal{C}$ tel que $A \subseteq C$.
Deuxi\`emement, le nombre de conteneurs est born\'e de mani\`ere logarithmique par les param\`etres du graphe, donnant $|\mathcal{C}| \le 2^{\epsilon N}$ pour un $\epsilon > 0$ donn\'e et $N$ suffisamment grand.
Troisi\`emement, le nombre d'ar\^etes induites par tout conteneur $C$ dans le graphe $G_N$ est strictement inf\'erieur \`a $\epsilon N^2$.
Parce que la densit\'e d'ar\^etes induites de chaque conteneur $C$ est \'evanescente, $C$ doit approximer \'etroitement l'un des ensembles ind\'ependants maximaux de $G_N$.
Les ensembles ind\'ependants maximaux de $G_N$ correspondant aux ensembles sans somme structurels sont l'ensemble des entiers impairs $\mathcal{O} = \{1, 3, 5, \dots, 2\lceil N/2 \rceil - 1\}$ et l'ensemble des grands entiers $\mathcal{L} = \{\lfloor N/2 \rfloor + 1, \dots, N\}$.
Par cons\'equent, tout conteneur $C \in \mathcal{C}$ satisfait $|C \setminus \mathcal{O}| < \delta N$ ou $|C \setminus \mathcal{L}| < \delta N$ pour un petit $\delta > 0$ qui d\'epend de $\epsilon$.
\end{proof}

\section{Architecture pour Autoformalisation}

Pour structurer ce probl\`eme dans Lean 4, nous d\'efinissons les types requis pour les ensembles et la propri\'et\'e d'\^etre sans somme.

\begin{lstlisting}[language=lean, basicstyle=\ttfamily\small, breaklines=true]
import Mathlib.Data.Finset.Basic
import Mathlib.Data.Nat.Basic

open Finset

-- Definition d'un ensemble sans somme
def IsSumFree (A : Finset Nat) : Prop :=
  \forall x \in A, \forall y \in A, (x + y) \notin A

-- L'enonce de la Conjecture de Cameron-Erdos
-- Nous definissons une fonction pour le nombre d'ensembles sans somme dans [1, N]
def numSumFreeSets (N : Nat) : Nat :=
  ((Icc 1 N).powerset.filter IsSumFree).card

-- L'enonce formel
theorem cameron_erdos_theorem :
  Exists (fun C : Real => \forall N : Nat, N > 0 -> (numSumFreeSets N : Real) \le C * 2^(N / 2)) := by
  admit

-- Lemme basique: L'ensemble des nombres impairs est sans somme
lemma odd_is_sum_free (N : Nat) :
  IsSumFree ((Icc 1 N).filter (fun x => x % 2 = 1)) := by
  admit
\end{lstlisting}

\end{document}
